% POR: Teixeira V. 
% VERSIÓN: 1.1
% ACTUALIZACIÓN: 24/03/2021
% -----------------------------------

\documentclass{article}[11pt]
\usepackage{indentfirst}
%\usepackage{fontspec}
\usepackage[spanish,es-tabla, es-ucroman]{babel}

\usepackage{calc}
\usepackage{float}
\usepackage{caption}
\usepackage{graphicx}
\usepackage{subfigure}
\usepackage{mathtools}
\usepackage{amsmath}
\usepackage{pgfplotstable} 
\usepackage{natbib}
\usepackage{url}
\usepackage{enumerate} 
\usepackage{amsmath}
\usepackage{graphicx}
\graphicspath{{figuras/}}
\usepackage{parskip}
\usepackage{fancyhdr}
\usepackage{vmargin}
\usepackage{changepage}
\usepackage{xcolor}
\usepackage{multicol}

\usepackage{siunitx}
\usepackage{hyperref}

% FUENTE
%setmainfont[BoldFont=arialbd.ttf,
%ItalicFont=ariali.ttf,
%BoldItalicFont=arialbi.ttf]{Arial.ttf}

% Estilo BIBLIOFGRAFIA
\bibliographystyle{abbrvnat}

% CONFIGURACIÓN DE LA HOJA
\setmarginsrb{4 cm}{2.5 cm}{2 cm}{2.5 cm}{1 cm}{1.5 cm}{1 cm}{0.5 cm}

% CONFIGURACIÓN DE SANGRÍA
\setlength{\parindent}{4mm}
\setlength{\parskip}{1mm}

%---------------------------------------------------------------------------
%	EN TODAS HOJAS
%---------------------------------------------------------------------------

\title{Nombre de la UC}				    	% Titulo
\author{Nombre del estudiante}	            % Nombre del estudiante				
\date{\today}						        % Fecha


\makeatletter
\let\thetitle\@title
\let\theauthor\@author
\let\thedate\@date
\makeatother

\pagestyle{fancy}
\fancyhf{}
\rhead{\theauthor}
\lhead{\thetitle}
\cfoot{\thepage}

%%%%%%%%%%%%%%%%%%%%%%%%%%%%%%%%%%%%%%%%%%%%%%%%%%%%%%%%%%%%%%%%%%%%%%%%%%%%%%%%%
%	                                CARÁTULA 
%%%%%%%%%%%%%%%%%%%%%%%%%%%%%%%%%%%%%%%%%%%%%%%%%%%%%%%%%%%%%%%%%%%%%%%%%%%%%%%%%

\pgfplotsset{compat=1.16}
\begin{document}

% CARÁTULA
\begin{titlepage}
\vspace*{-3 cm}

\begin{adjustwidth}{2.5cm}{}    % aumenta sangría sin cambiar margenes
\includegraphics[scale = 0.2]{./figuras/logo_itr.png}    % logo utec
\end{adjustwidth}
\vspace{1 cm}

%%%% TITULO DEL TRABAJO
\begin{adjustwidth}{3cm}{}

 \definecolor{Micolor1}{RGB}{0,176,240}                     % color
 \noindent \rule{11.9cm}{0.4mm} 

 \begin{flushleft}                               %alineado a ezquierda
 
 \Huge \bfseries \textcolor{Micolor1}{MODELO DE INFORME - Carrera de Ingeniería en Energías Renovables}\\[0.7 cm]
 
 \end{flushleft}
 
 \noindent \rule{12cm}{0.4mm} \\ [0.1 cm] 
 \definecolor{Micolor1}{RGB}{0,176, 240}
 \textbf{\large \textcolor{Micolor1}{SEDE:}} \large{DURAZNO}

\end{adjustwidth}
\vspace{2 cm}


%%%% NOMBRE DE LA UC
\begin{adjustwidth}{3cm}{}
  \textmd {\Large Nombre de la UC:}\\[0.3 cm]
  \includegraphics[scale = 0.4]{./figuras/logo_tec.png}	 
\vspace{1 cm}


%%%% NOMBRE DE AUTOR Y TUTOR
\begin{multicols}{2}     

%NOMBRE DE AUTOR
\definecolor{Micolor1}{RGB}{0,176, 240}
\setlength{\parindent}{0cm} \textbf{ \large \textcolor{Micolor1}{Estudiante:}} \large{Nombre y apellido}

%NOMBRE TUTOR
\columnbreak
\textbf{ \large \textcolor{Micolor1}{Docentes:}} \large{XXXXX}\\
\columnbreak
\hspace{2.2cm} YYYYY
\end{multicols}

\vspace{2 cm}
	      \end{adjustwidth}

% FECHA
\begin{center}
\definecolor{Micolor1}{RGB}{0,176, 240}
\textbf{\small \textcolor{Micolor1}{Fecha:}} \small{\thedate}

	\vfill
\end{center}

\end{titlepage}

%%%%%%%%%%%%%%%%%%%%%%%%%%%%%%%%%%%%%%%%%%%%%%%%%%%%%%%%%%%%%%%%%%%%%%%%%%%%%%%%%
%	                     INICIO DEL DOCUMENTO  
%%%%%%%%%%%%%%%%%%%%%%%%%%%%%%%%%%%%%%%%%%%%%%%%%%%%%%%%%%%%%%%%%%%%%%%%%%%%%%%%%

%%%%%%% RESUMEN
\renewcommand{\thepage}{\roman{page}}           %numero de pag romano

\begin{center}
\normalsize \textbf{Resumen}
\end{center}

\noindent Deberá presentar el objetivo del trabajo, con una breve formulación de la metodología utilizada y las conclusiones o resultados obtenidos. 
\vspace{0.3cm}

\noindent \underline{Palabras clave:} Citar tres palabras que representa el estudio.
\newpage

%%%%%%% INDICE DE FIGURAS
\listoffigures
\addcontentsline{toc}{section}{Índice de figuras}
\newpage
\thispagestyle{empty} % quita numero de la pagina

%%%%%%% ÍNDICE GENERAL
\tableofcontents
\newpage
\pagenumbering{arabic} % cera numero de la pagina

%%%%%%%%%%%%%% 
% SECCIÓN 1
%%%%%%%%%%%%%%
\section{Características del índice}

Deberá tener las siguientes características:\\
\begin{itemize}
\item Constar de un cuadro esquemático en el cual se detallen las divisiones que componen el trabajo, remitiendo a la página donde aparece desarrollado el contenido.
\item Estar enumerados todos los capítulos, secciones o partes del trabajo, y la lista de referencias bibliográficas. Incluye los anexos, listas de cuadros que se hubieren incorporado.
\item Ser un fiel reflejo del cuerpo de la obra.
\end{itemize}

%%%%%%%%%%%%%% 
% SECCIÓN 2
%%%%%%%%%%%%%%
\section{Introducción}

A continuación, se destaca algunos criterios importantes para redactar la introducción:\\
\begin{itemize}
\item Presentar el contexto del tema, la problemática y su importancia.
\item Redactar de forma precisa la información que se ha tenido en cuenta para el desarrollo del trabajo, respondiendo preguntas básicas sobre el tema (ej. \textit{¿qué es?, ¿para qué sirve?, ¿cuál es la importancia?}, y otras).
\end{itemize}

%%%%%%%%%%%%%% 
% SECCIÓN 3
%%%%%%%%%%%%%%
\section{Objetivos}
Los objetivos fijarán los alcances del problema, indicando lo que se pretende lograr al finalizar el trabajo. Antes de plantear los objetivos es importante saber que los mismos:\\

\begin{itemize}
\item Deben ser concretos y viables.
\item Pueden dividirse en generales y específicos. Los específicos están relacionados con los problemas derivados, o constituir aspectos parciales del objetivo general.
\end{itemize}

%%%%%%%%%%%%%% 
% SECCIÓN 4
%%%%%%%%%%%%%%
\section{Marco teórico}
\begin{itemize}
\item Incluir las teorías, postulados, conceptos, y otros aspectos que sostienen el abordaje del objeto de estudio y toda la investigación en general. 
\item Presentar antecedentes del tema, haciendo referencias a estudios relevantes, presentando de manera resumida la metodología utilizada y los resultados principales obtenidos. 
\item Está basado principalmente en distintas referencias bibliográficas. 
\end{itemize}

%%%%%%%%%%%%%% 
% SECCIÓN 5
%%%%%%%%%%%%%%
\section{Metodología}
Esta sección debe contestar preguntas como por ejemplo: \textit{¿Cómo?, ¿Donde?, ¿Con qué? y ¿Cuanto?}.\\
\begin{itemize}
\item En estudios experimentales presentar los procedimiento de armado, mecanismos de obtención de datos, herramientas de análisis de la información generada.
\item En estudios de caso explicar cómo se obtiene y analiza los aspectos variados del caso.
\item Presentar los indicadores que serán utilizados para la análisis de los resultados.
\item Describir los equipos, materiales y/o software utilizados.
\end{itemize}

%%%%%%%%%%%%%% 
% SECCIÓN 6
%%%%%%%%%%%%%%
\section{Resultados y discusión}
Los resultados obtenidos deberán ser presentados de forma clara, sintética y precisa. \\
\begin{itemize}
\item La presentación de los resultados puede ser textual, sin embargo siempre que posible es recomendable la utilización de tablas, gráficas, figuras, dibujos acompañados de la descripción correspondiente.
\item La discusión de los resultados debe ser realizada a partir de la análisis de resultados, acompañada de su fundamentación teórica. Se esperan comentarios a cerca de los aportes realizados a partir los resultados obtenidos en el trabajo. Además se recomienda relacionar con estudios similares en la literatura.
\item Evaluar si los resultados obtenidos corresponden a los objetivos planteados.
\end{itemize} 

%%%%%%% 1 SUBSECCIÓN %%%%%%
\subsection{Sobre las figuras}
Las figuras deben cumplir con los siguientes criterios: \\

\begin{itemize}
\item Estar numeradas y con título en la parte inferior del texto.
\item Deben ser explicadas y citadas con su respectivo numero antes de su presentación en el documento. 
\item Tener buena resolución. En el caso en que hayan texto, este debe estar en el idioma local.
\item Ser citadas con su respectivo número antes de su presentación en el documento. 
\item Las gráficas deben tener los ejes con título, unidades de medida y leyendas si necesario. El tamaño del título de los ejes y de la leyenda deben ser visibles al lector.\\
\end{itemize} 

La Figura \ref{fig:logo} muestra un ejemplo.

\begin{figure}[!ht]
\centering
\includegraphics[width=0.2\textwidth]{figuras/Isotipo_NEGRO.png}
\caption{Logo UTEC}
\label{fig:logo}
\end{figure}

%%%%%%% 2 SUBSECCIÓN %%%%%%

\subsection{Sobre las tablas}
Sigue algunos criterios a tener en cuenta en el momento de presentar tablas:\\

\begin{itemize}
\item Deben estar numeradas y con título en la parte inferior del texto.
\item Deben ser explicadas y citadas con su respectivo número antes de su presentación en el documento. 
\item Las magnitudes de estudio deben estar identificadas (ej. datos, sitio, potencia, y otros), acompañada de su unidad de medida siempre que haya. En caso de que se utilicen símbolos, estos deben ser explicados en el texto. 
\item Los valores  deben ser escritos teniendo en cuenta las cifras significativas que indique el docente responsable.
\item Para números muy grandes o pequeños se debe usar notación científica.\\
\end{itemize}

Ejemplos:

\begin{table}[htb]
\begin{center}
\renewcommand\arraystretch{1.3}
\centering
\begin{tabular} {c c c c}
 \hline
N° & t (s) & I (W/m$^{2}$) & d (m)\\
 \hline
1 & 1,23 & 4200 & 1,20\\
2 & 1,25 & 3997 & 1,50\\
3 & 1,23 & 4200 & 1,20\\
 \hline
 \end{tabular}
 \caption{Valores experimentales de la práctica X}
\label{ta:tab1}
    \end{center}
 \end{table}
 
\begin{table}[htb]
\begin{center}
\begin{tabular}{l c c}
\hline
Mes & $\mathrm{HSP_{PH} (h)}$ & $\mathrm{HSP_{PI} (hrs)}$ \\
\hline
Enero & 224 & 213\\ 
Julio & 224 & 213\\ 
Diciembre & 224 & 213\\ 
\hline
\end{tabular}
\caption{Ejemplo 2 de tabla}
    \end{center}
\label{tab:ValoresHSP}
\end{table}

%%%%%%%%%%%%%% 
% SECCIÓN 7
%%%%%%%%%%%%%%
\section{Conclusiones}
Este sección debe cerrar la propuesta del trabajo. Algunos criterios para la redacción: \\

\begin{itemize}
\item Ser objetivo. 
\item Retomar la propuesta inicial y vincular con los resultados obtenidos.
\item Realizar un análisis critico, proponer mejoras y estudios futuros.
\end{itemize}

%%%%%%%%%%%%%% 
% SECCIÓN 8
%%%%%%%%%%%%%%
\section{Características de la bibliografía}
La sección bibliográfica debe estar ordenada alfabéticamente por apellido del autor. Se recomienda la utilización del formato APA (American Psychological Association).

%%%%%%%%%%%%%% 
% SECCIÓN 9
%%%%%%%%%%%%%%
\section{Características de los anexos}

\begin{itemize}
\item Pueden ser documentos de distinta índole, que aporten información complementaria o aclaratoria de cierto aspecto de la obra. 
\item Si fueron extraídos de fuentes consultadas, se deberán incluir sus referencias. 
\item Deben estar incorporados en el orden en que fueron mencionados en el desarrollo del trabajo.
\end{itemize}


\newpage
\bibliography{Biblio}
\nocite{*}
\end{document}
